\documentclass[10pt,a4paper]{amsart}
\usepackage{amsmath,amssymb,amsthm,graphicx}
\usepackage{verbatim}
\newenvironment{reference}{\comment}{\endcomment}
%\newenvironment{reference}{}{}
\newenvironment{slogan}{\comment}{\endcomment}
\newenvironment{history}{\comment}{\endcomment}

% For commutative diagrams we use Xy-pic
\usepackage[all]{xy}

% We use 2cell for 2-commutative diagrams.
\xyoption{2cell}
% Theorem environments.
%
\theoremstyle{plain}
\newtheorem{theorem}[subsection]{Theorem}
\newtheorem{proposition}[subsection]{Proposition}
\newtheorem{lemma}[subsection]{Lemma}

\theoremstyle{definition}
\newtheorem{definition}[subsection]{Definition}
\newtheorem{example}[subsection]{Example}
\newtheorem{exercise}[subsection]{Exercise}
\newtheorem{situation}[subsection]{Situation}

\theoremstyle{remark}
\newtheorem{remark}[subsection]{Remark}
\newtheorem{remarks}[subsection]{Remarks}

\numberwithin{equation}{subsection}

% Macros
%
\def\lim{\mathop{\mathrm{lim}}\nolimits}
\def\colim{\mathop{\mathrm{colim}}\nolimits}
\def\Spec{\mathop{\mathrm{Spec}}}
\def\Hom{\mathop{\mathrm{Hom}}\nolimits}
\def\Ext{\mathop{\mathrm{Ext}}\nolimits}
\def\SheafHom{\mathop{\mathcal{H}\!\mathit{om}}\nolimits}
\def\SheafExt{\mathop{\mathcal{E}\!\mathit{xt}}\nolimits}
\def\Sch{\mathit{Sch}}
\def\Mor{\mathop{\mathrm{Mor}}\nolimits}
\def\Ob{\mathop{\mathrm{Ob}}\nolimits}
\def\Sh{\mathop{\mathit{Sh}}\nolimits}
\def\NL{\mathop{N\!L}\nolimits}
\def\CH{\mathop{\mathrm{CH}}\nolimits}
\def\proetale{{pro\text{-}\acute{e}tale}}
\def\etale{{\acute{e}tale}}
\def\QCoh{\mathit{QCoh}}
\def\Ker{\mathop{\mathrm{Ker}}}
\def\Im{\mathop{\mathrm{Im}}}
\def\Coker{\mathop{\mathrm{Coker}}}
\def\Coim{\mathop{\mathrm{Coim}}}

% Boxtimes
%
\DeclareMathSymbol{\boxtimes}{\mathbin}{AMSa}{"02}

%
% Macros for moduli stacks/spaces
%
\def\QCohstack{\mathcal{QC}\!\mathit{oh}}
\def\Cohstack{\mathcal{C}\!\mathit{oh}}
\def\Spacesstack{\mathcal{S}\!\mathit{paces}}
\def\Quotfunctor{\mathrm{Quot}}
\def\Hilbfunctor{\mathrm{Hilb}}
\def\Curvesstack{\mathcal{C}\!\mathit{urves}}
\def\Polarizedstack{\mathcal{P}\!\mathit{olarized}}
\def\Complexesstack{\mathcal{C}\!\mathit{omplexes}}
% \Pic is the operator that assigns to X its picard group, usage \Pic(X)
% \Picardstack_{X/B} denotes the Picard stack of X over B
% \Picardfunctor_{X/B} denotes the Picard functor of X over B
\def\Pic{\mathop{\mathrm{Pic}}\nolimits}
\def\Picardstack{\mathcal{P}\!\mathit{ic}}
\def\Picardfunctor{\mathrm{Pic}}
\def\Deformationcategory{\mathcal{D}\!\mathit{ef}}
\usepackage{cleveref}
\crefname{originallemma}{original lemma}{original lemmas}
\Crefname{originallemma}{Original lemma}{Original lemmas}
\makeatletter
\newlabel{chow-lemma-intersect-properly}{{0FFD}{source}{}{sourcepointer-1}{}}
\newlabel{chow-lemma-intersect-properly@cref}{{[originallemma][0][]0FFD}{[1][0][]source}{}{}{}}
\makeatother
\crefname{originalreference}{original reference}{original references}
\Crefname{originalreference}{Original reference}{Original references}
\makeatletter
\newlabel{lemma-category-correspondences}{{0FG2}{source}{}{sourcepointer-2}{}}
\newlabel{lemma-category-correspondences@cref}{{[originalreference][0][]0FG2}{[1][0][]source}{}{}{}}
\makeatother
\title{The contravariant functor from smooth projective schemes to correspondences}
\author{The Stacks Project Authors}
\date{stacks-project:weil.tex:local-snapshot-20261012; GNU FDL 1.2 or later}
\begin{document}
\maketitle
\noindent\textit{Editorial note.} All delivered mathematical text is the author's own, copied from the pinned Stacks Project chapter file weil.tex; no statement or proof step was rewritten, reordered or removed.\par\medskip
\noindent\textit{Editorial note.} The counted proof is the author's complete argument: the behaviour on identities, the reduction of functoriality to an intersection product, and the identification of the triple graph.\par\medskip
\noindent\textit{Editorial note.} The label lemma-contravariant-functor names the counted statement, and the proof environment immediately following it in the source is delivered verbatim as its proof.\par\medskip
\noindent\textit{Editorial note.} The counted statement is carried by the author's own lemma environment in the chapter source.\par\medskip
\noindent\textit{Editorial note.} Pre-open gate: python operations/agent-authoring/tools/precard\_gate.py --source weil.tex --label lemma-contravariant-functor --start 609 --end 642 --local-id 28145 returned RESULT=OK (L28145 free, occupancy FREE, public interval clear). Pre-assembly tag self-check: canonical tag resolved and not present in reference/public-items.jsonl.\par\medskip
\noindent\textit{Editorial note.} The statement and proof are copied from the pinned chapter file; the cited inputs are kept as external pointers with their canonical Stacks tags.\par\medskip
\noindent\textit{Allowed standard prerequisites (editorial).}\begin{itemize}
\item Identities are sent to identities. Same chapter, the proof of the construction of the category of correspondences, where the behaviour on identities is verified. Provides the identity part of the functoriality.
\item The intersection product of the graph with a product graph. Chow Homology, the lemma on intersecting cycles that meet properly. Identifies the class of the image of the triple graph with the intersection product of the two pulled back graphs.
\end{itemize}\medskip
\ifdefined\hypertarget\hypertarget{sourcepointer-1}{}\fi
\noindent\textit{Source reference.} Original lemma 0FFD: sources/stacks-chow/tags.txt, source lines 19283--19283. Prerequisite pointer invoked at proof line(s) 635; it names a standard fact of a sibling chapter that the counted argument uses, and every step of the counted proof built on it is retained verbatim in the delivered body.\par\medskip
\ifdefined\hypertarget\hypertarget{sourcepointer-2}{}\fi
\noindent\textit{Source reference.} Original reference 0FG2: sources/stacks-weil/pinned-original-source.tex, source lines 550--550. Prerequisite pointer invoked at proof line(s) 618; it names a standard fact of the same chapter that the counted argument uses, and every step of the counted proof built on it is retained verbatim in the delivered body.\par\medskip
\begin{lemma}
\label{lemma-contravariant-functor}
There is a contravariant functor from the category of smooth
projective schemes over $k$ to the category of correspondences
which is the identity on objects and sends $f : Y \to X$ to
the element $[\Gamma_f] \in \text{Corr}^0(X, Y)$.
\end{lemma}

\begin{proof}
In the proof of Lemma \ref{lemma-category-correspondences}
we have seen that this construction sends identities to
identities. To finish the proof we have to show if $g : Z \to Y$
is another morphism of smooth projective schemes over $k$, then we have
$[\Gamma_g] \circ [\Gamma_f] = [\Gamma_{f \circ g}]$ in
$\text{Corr}^0(X, Z)$. Arguing as in the proof of
Lemma \ref{lemma-category-correspondences} we see that it
suffices to show
$$
[\Gamma_{f \circ g}] =
\text{pr}_{13, *}([\Gamma_f \times Z] \cdot [X \times \Gamma_g])
$$
in $\CH^*(X \times Z)$ when $X$, $Y$, $Z$ are integral.
Denote $Z' \subset X \times Y \times Z$ the image of the closed immersion
$(f \circ g, g, 1) : Z \to X \times Y \times Z$.
Then $Z' = \Gamma_f \times Z \cap X \times \Gamma_g$
scheme theoretically and we conclude using
Chow Homology, Lemma \ref{chow-lemma-intersect-properly}
that
$$
[Z'] = [\Gamma_f \times Z] \cdot [X \times \Gamma_g]
$$
Since it is clear that $\text{pr}_{13, *}([Z']) = [\Gamma_{f \circ g}]$
the proof is complete.
\end{proof}

\end{document}
